\begin{answers}{ch:math}{数学の準備：記号・微分・確率・情報量}
\item 内積は 4、ノルムは $\norm{\va}=3$、$\norm{\vb}=\sqrt{5}\approx 2.236$、コサイン類似度は約 0.596 です。
内積は式\eqref{eq:math-inner}より $1\times 2+2\times 0+2\times 1=4$ です。
ノルムは式\eqref{eq:math-norm}より $\sqrt{1+4+4}=3$ と $\sqrt{4+0+1}=\sqrt{5}$ です。
式\eqref{eq:math-cosine}に代入すると $4/(3\sqrt{5})\approx 0.596$ です。
正なので、二つの矢印は、直角よりは同じ向きに近いと言えます。
\item 勾配は $\nabla f=(2xy,\ x^{2}+3y^{2})^\T$ で、点 $(1,2)$ では $(4,13)^\T$ です。
$x$ で偏微分するときは $y$ を定数と見るので、$x^{2}y$ は $2xy$ に、$y^{3}$ は 0 になります。
$y$ で偏微分するときは $x$ を定数と見るので、$x^{2}y$ は $x^{2}$ に、$y^{3}$ は $3y^{2}$ になります。
$(1,2)$ を代入すると、$2\times1\times2=4$ と $1+3\times 4=13$ です。
\file{ch01/numgrad.py} の \code{f} を \code{return x**2 * y + y**3} に変え、\code{numerical_grad(f, 1.0, 2.0)} とすると、\code{[4. 13.]} が出て一致しました。
\item $H(p)=1.5\log 2\approx 1.040$、$H(p,q)=1.75\log 2\approx 1.213$、$\KL{p}{q}=0.25\log 2\approx 0.173$ です。
$\log 0.5=-\log 2$、$\log 0.25=-2\log 2$ を使うと、手で計算できます。
$H(p)=0.5\log 2+0.25\times 2\log 2+0.25\times 2\log 2=1.5\log 2$ です。
$H(p,q)=0.5\times 2\log 2+0.25\log 2+0.25\times 2\log 2=1.75\log 2$ です。
KL は、差 $H(p,q)-H(p)=0.25\log 2$ です（式\eqref{eq:math-kl}）。
\file{ch01/entropy.py} の関数でも、同じ値が出ました。
なお、この例では $\KL{q}{p}$ も 0.173 になります。$p$ と $q$ が成分を入れ替えただけの関係だからで、一般には二つは等しくありません。
\item 「ない。」の確率は、もとの 0.559 から、2 倍で 0.934 に上がり、半分で 0.055 に下がりました。
\file{ch01/softmax.py} の \code{z = W @ h} の次に \code{z = 2 * z}（または \code{z = 0.5 * z}）を加えて実行した結果です。
確率の比は $e^{z_i-z_j}$ なので（\secref{sec:math-softmax}）、ロジットを 2 倍すると比は 2 乗に、半分にすると平方根になります。
「ない。」と「ないから」の比 4.62 は、2 倍では約 21 に、半分では約 2.1 になりました。
2 倍では分布がとがり、予測のエントロピーは 2.204 から 0.318 に下がります。
半分では、4096 個のトークンに確率が広く散らばり、エントロピーは 7.213 に上がります。
いずれの場合も、確率の大きい順は変わりません。
2 倍は\chapref{ch:lm}の温度 $\tau=0.5$、半分は $\tau=2$ に当たります（式\eqref{eq:lm-temperature}）。
\item 損失は $-\log 0$ で、限りなく大きく（無限大に）なります。
$-\log q$ は、$q$ を 0 に近づけるといくらでも大きくなるからです（\figref{fig:math-info}）。
たとえば $q=10^{-10}$ でも、損失は約 23 です。
そのため、学習は「起こりうるトークンには、少しでも確率を残しておく」ようにモデルを押します。
確信のない場面で一つに決めつけると、外れたときの罰がとても大きいからです。
なお、softmax の出力は必ず正なので（\figref{fig:math-softmax}）、実際のモデルでは確率がちょうど 0 になることはありません。
\end{answers}
